\chapter{Aboriginal and Torres Strait Islander Social and Emotional Wellbeing}

\section*{Definition and Overview}
Social and emotional wellbeing (SEWB) is the Aboriginal and Torres Strait Islander concept of mental health and wellbeing. It is holistic: wellbeing is understood to arise from connections between body, mind, and emotions, and from relationships with family, community, culture, Country, spirit, and ancestors. This framework is broader than the Western model of mental health, in which distress is often seen primarily as an individual illness. Within the SEWB framework, difficulties in one area of life affect the whole person, and healing involves restoring balance across all these connections. Care that supports SEWB therefore respects these relationships and works with families and communities, not with the individual in isolation.

\section*{Epidemiology}
Aboriginal and Torres Strait Islander peoples experience higher rates of psychological distress and mental illness than the wider population, and rates of self-harm and suicide are substantially higher, particularly among young people. Around one-third of Aboriginal and Torres Strait Islander adults report high or very high levels of psychological distress. Suicide is a leading cause of premature death in these communities, and its effects cascade through families and communities. At the same time, connection to culture, family, community, and Country is strongly protective, and many people report high levels of wellbeing, demonstrating that adversity and strength can coexist.

\section*{Aetiology and Pathophysiology}
The drivers of distress in Aboriginal and Torres Strait Islander communities are primarily social, cultural, and historical.

\begin{itemize}
    \item \textbf{Social determinants:} housing, income, education, employment, and access to services
    \item \textbf{Intergenerational trauma:} the enduring effects of colonisation, dispossession, and the forced removal of children, transmitted across generations
    \item \textbf{Grief and loss:} ongoing and unresolved losses of Country, family, language, and culture
    \item \textbf{Racism and discrimination:} experienced in daily life, workplaces, and services
    \item \textbf{Cultural disconnection:} reduced opportunity to practise culture, language, and connection to Country
    \item \textbf{Family and community stress:} including family violence, substance use, and heavy caring responsibilities
    \item \textbf{Physical health problems:} chronic illness and poor physical health interact with emotional wellbeing in both directions
\end{itemize}

\section*{Clinical Features}
Distress may be experienced and expressed in many ways, and it is important to recognise that strong connection to family and community can coexist with significant personal distress.

\begin{itemize}
    \item Physical symptoms such as headaches, stomach aches, and fatigue with no clear physical cause
    \item Low mood, anxiety, irritability, or a sense of feeling flat or empty
    \item Changes in sleep, appetite, or energy levels
    \item Withdrawal from family, community, or cultural activities
    \item Increased use of alcohol, tobacco, or other drugs
    \item Difficulty concentrating, or feeling disconnected or numb
    \item Thoughts of self-harm or suicide, which must always be taken seriously
    \item Conversely, many people maintain strong wellbeing through family, community, and cultural connection
\end{itemize}

\section*{Differential Diagnosis}
A range of explanations should be considered in any presentation.

\begin{itemize}
    \item Physical illness that may be causing the symptoms
    \item Depression, anxiety, or another diagnosable mental health condition
    \item Grief and loss, which may be ongoing and cumulative rather than a single event
    \item The effects of alcohol or other drugs
    \item Family violence or other serious stressors at home
    \item Cultural disconnection, which may be driving distress without a mental illness being present
    \item A combination of several of these, which is common
\end{itemize}

\section*{When to Seek Medical Care}
Help should be sought if symptoms persist, interfere with daily life, or raise any concern about safety. Help is available through a GP, an Aboriginal Community Controlled Health Organisation, or a dedicated social and emotional wellbeing service, and seeking it early is strongly encouraged.

\textbf{Call 000 (or local emergency services) for:}
\begin{itemize}
    \item Thoughts of suicide or plans to self-harm
    \item A person who has harmed themselves or is in immediate danger
    \item Severe agitation, confusion, or an inability to keep the person safe
    \item Any mental health crisis in which safety is a concern
\end{itemize}

For someone in crisis, stay with them, remove any means of harm, and seek urgent professional help.

\section*{Diagnosis and Investigations}
Assessment is built on trust, listening, and a culturally safe approach.

\begin{itemize}
    \item \textbf{Culturally safe assessment:} taking time to build trust, often using yarning and conversation rather than formal questioning
    \item \textbf{Listening to the person's own account} of their connections, strengths, and what matters to them
    \item \textbf{Involving family and community} in the assessment where the person wishes
    \item \textbf{Screening for distress and risk,} including asking directly and without judgement about thoughts of self-harm or suicide
    \item \textbf{Considering both SEWB and mental health frameworks,} since both can be useful in planning care
    \item \textbf{Physical health checks} where symptoms may have a physical cause
\end{itemize}

\section*{Management}
Management is most effective when it is holistic, community-based, and delivered by people who understand culture and Country.

\begin{itemize}
    \item \textbf{Social and emotional wellbeing services,} often staffed by Aboriginal health workers and Aboriginal counsellors
    \item \textbf{Culturally safe psychological therapy,} adapted so that it aligns with the person's values, family, and community
    \item \textbf{Trauma-informed care} that recognises the effects of past and intergenerational trauma
    \item \textbf{Connection to culture and Country:} cultural activities, men's and women's groups, and on-Country healing programs
    \item \textbf{Family involvement} in care and in the healing journey
    \item \textbf{Support for underlying social issues} such as housing, employment, finances, and violence
    \item \textbf{Coordinated care} between mainstream and Aboriginal community-controlled services
    \item \textbf{Regular follow-up,} especially in the period after a crisis
\end{itemize}

\section*{Complications}
\begin{itemize}
    \item Worsening depression, anxiety, and other mental health conditions when support is delayed
    \item Alcohol and other drug problems, which often coexist with distress
    \item Self-harm and suicide, which devastate families and communities
    \item Physical health problems linked to chronic stress and psychological distress
    \item Breakdown of family and community relationships
    \item Stigma and shame, which prevent people from seeking help early
\end{itemize}

\section*{Prognosis and Outcomes}
Many people recover well from episodes of distress, particularly when they receive timely, culturally safe support and retain strong family and community connection. Recovery is more likely when care addresses the whole person, including culture, Country, and the social causes of stress, rather than symptoms alone. Early help improves outcomes, and people who have been through a crisis can go on to live full and connected lives. Healing within an SEWB framework is often understood as a journey that involves family and community as much as the individual.

\section*{Prevention and Self-Care}
\begin{itemize}
    \item Strengthen connection to culture, family, community, and Country
    \item Take part in activities that bring purpose, meaning, and belonging
    \item Get enough sleep, eat well, and stay physically active
    \item Limit alcohol and other drugs, which can worsen distress
    \item Talk to trusted family members, elders, or community members when things are difficult
    \item Seek help early, before distress builds up
    \item For communities: support programs that strengthen culture, address social disadvantage, and train local people as healers and workers
\end{itemize}

\section*{Sources and Further Reading}
The following organisations provide further information on social and emotional wellbeing and Aboriginal and Torres Strait Islander mental health.

\begin{itemize}
    \item Australian Institute of Health and Welfare. \textit{Social and emotional wellbeing of Aboriginal and Torres Strait Islander people.} https://www.aihw.gov.au/
    \item Australian Government Department of Health. \textit{Indigenous mental health and suicide prevention programs.} https://www.health.gov.au/
    \item NACCHO. \textit{Social and emotional wellbeing in Aboriginal community-controlled health.} https://www.naccho.org.au/
    \item Beyond Blue. \textit{Mental health resources for Aboriginal and Torres Strait Islander peoples.} https://www.beyondblue.org.au/
    \item Headspace. \textit{Support for young Aboriginal and Torres Strait Islander people.} https://headspace.org.au/
    \item World Health Organization. \textit{Mental health and social determinants.} https://www.who.int/
\end{itemize}
